\begin{claim}{C.42}{Proved}{The actual boundary quotient has proper line components upstairs}
Let $G$ be generated by a based plane loop as in C.28,
let $O$ be an open $G$-orbit, and put
$L=\overline G\leq\operatorname{Homeo}(\R^2)$.
Fix $x\in O$ and $J=L_x$. For
$z\in\operatorname{Fix}(J)$ whose $G$-orbit $B=Gz$
is a nontrivial intrinsic line or a circle, the formula
\[
 \Phi_z(lx)=lz
\]
defines the actual continuous open quotient
$O=L/J\to B=L/L_z$, with the intrinsic topology on $B$.
For the universal coverings $p:\widetilde O\to O$
and $q:\R\to B$, every lift
$F:\widetilde O\to\R$ of $\Phi_zp$ is open and onto,
and every connected component of every fiber of $F$
is a properly embedded line in the covering plane
$\widetilde O$.
Every nondegenerate compact connected component of
a fiber of $\Phi_z$ is a circle.
Whole fibers of $\Phi_z$ can be disconnected, including
for quotients obtained as limits of the original
commuting operators.
\end{claim}

\begin{proof}
\emph{The quotient and its topology.}
Every element of $L$ preserves $O$: forward and inverse
limits preserve its boundary, while a forward limit
sends an interior point into $\overline O$ and cannot
send it to that boundary. The Polish orbit-map theorem
used in C.26 and C.30 therefore makes $l\mapsto lx$
an open quotient $L\to O$.
If $l_1x=l_2x$, then $l_2^{-1}l_1\in J$ fixes $z$,
so $\Phi_z$ is well defined and continuous.
More precisely, the continuous action map
\[
 \operatorname{Fix}(J)\times L\longrightarrow\R^2,
 \qquad (z,l)\longmapsto lz
\]
descends through the open quotient
$\operatorname{id}\times(L\to O)$.
Consequently
\[
 \operatorname{Map}_L(O,\R^2)
       \xrightarrow[\cong]{\,\operatorname{ev}_x\,}
       \operatorname{Fix}(J)
                                                        \tag{C.42.1}
\]
is a homeomorphism in the compact-open topology:
equivariance proves bijectivity, and the displayed
joint action proves continuity of the inverse.
For $l\in L$, choose $g\in G$ with $gx=lx$.
Then $g^{-1}l\in J$, hence $lz=gz$.
Thus the range is exactly $Lz=Gz=B$.

The map into an intrinsic line is continuous in its
intrinsic topology. Indeed, the image of a compact
coordinate disk in $O$ is a compact Peano continuum
in $B$. By C.28 it contains neither a Jordan curve
nor a triod, and hence is a point or a compact arc;
on such an arc the intrinsic and ambient topologies
agree. These disks give neighborhoods at every point
of $O$. For a circle the two topologies already agree.
Each $l\in L$ restricts to an intrinsic homeomorphism
of $B$, because compact arcs are carried to compact
arcs. For $b=gz$ one has
\[
                  lb=\Phi_z(lgx).
                                                        \tag{C.42.2}
\]
This makes every point evaluation of the restriction
$L\to\operatorname{Homeo}(B)$ continuous.
Finite ordered partitions of compact intervals show
that pointwise convergence of monotone homeomorphisms
to a homeomorphism is uniform on compact intervals;
apply the same argument to inverses. Finite cyclic
partitions give the circle assertion. Thus restriction
is continuous in the compact-open group topology.
All these restrictions preserve orientation, since
the original loop does and the restriction map is
continuous.

Apply the Polish orbit-map theorem to this transitive
action on the Polish one-manifold $B$. It identifies
$B$ with $L/L_z$ in its stated topology. Since
$J\subseteq L_z$, the map $\Phi_z$ is the open quotient
$lJ\mapsto lL_z$.
In particular, if operators $h_n$ commute with $G$
and $h_nx\to z$, then they commute with $L$ and
(C.42.1) gives $h_n|_O\to\Phi_z$ uniformly on compact
subsets of $O$.

\emph{Both covering maps and the compatible group.}
The universal cover of the connected open planar
surface $O$ is a plane. Fix a lift
\[
 F:\widetilde O\longrightarrow\R,\qquad
                      qF=\Phi_zp.
                                                        \tag{C.42.3}
\]
For a line $B$, take $q$ to be an identification.
Local covering charts and openness of $\Phi_zp$
show that $F$ is open.
Consider the following subgroup, with the topology
induced from the indicated product:
\[
\begin{split}
\mathcal L=\{(l,a,b)\in
 L\times\operatorname{Homeo}(\widetilde O)
       \times\operatorname{Homeo}^+(\R):\quad
 pa&=(l|_O)p,\\
 qb&=(l|_B)q,\qquad Fa=bF\}.
\end{split}
                                                        \tag{C.42.4}
\]
The three equations are closed conditions, as may
be checked by evaluation on countable dense sets.
Hence $\mathcal L$ is a Polish group, and its
actions by $a$ and by $b$ are continuous.
Every lift $a$ of $l|_O$ has a unique compatible
lift $b$ of $l|_B$: the maps $Fa$ and $bF$ lift
the same map to $B$, and the target deck translation
is chosen to make them agree at one point.
Uniqueness of lifts then gives equality everywhere.
For a line target there is no target deck ambiguity.

The action of $\mathcal L$ on $\widetilde O$ is
transitive. Choose $l$ carrying the projections of
two specified points to one another and adjust a
source lift by a deck transformation to carry the
specified points themselves.
The image of $F$ is a nonempty open interval $I$.
It is invariant under every compatible $b$.
If $B$ is a circle and $I$ has a finite endpoint,
every increasing $b$ fixes that endpoint, so its
image under $q$ would be fixed by every element
of the transitive $L$-action on $B$. This is
impossible. Therefore $I=\R$; for a line target
surjectivity is already immediate from (C.42.3).

\emph{A component that separates the plane.}
Fix $t\in\R$ and write $E=F^{-1}(t)$.
Choose $u,v$ with $F(u)<t<F(v)$. The closed set
$E$ separates $u$ from $v$. Frankel's planar
separation lemma \cite[Lemma 2.7, p.~1258]{C_Frankel2015}
gives a connected component $C$ of $E$ separating
them. The lemma applies to unbounded closed sets.
For completeness, take the complementary domain
$U$ of $E$ containing $u$, and then the component
$V$ of $\widetilde O\setminus\overline U$ containing
$v$. The planar Brouwer property makes
$\partial V$ connected. It lies in $E$ and
separates $u,v$, so the component of $E$ containing
$\partial V$ has the required property.

Let $\mathcal K$ be the closed subgroup of
$\mathcal L$ whose $b$-coordinate fixes $t$.
It acts transitively on $E$: transitivity on the
covering plane finds an element carrying any
two points of $E$ to one another, and
$Fa=bF$ forces $b(t)=t$.
Thus all components of $E$ are ambiently
homeomorphic to $C$, and all separate the plane.
The stabilizer $\mathcal K_C$ is a closed Polish
subgroup and acts transitively on $C$.
Closedness follows by applying forward and inverse
limits to the closed set $C$.

\emph{Accessibility from each fixed complementary domain.}
Let $D$ be any component of
$\widetilde O\setminus C$.
Its stabilizer in $\mathcal K_C$ is open:
for a fixed $d\in D$, every sufficiently small
group element sends $d$ into $D$ and therefore
sends the whole complementary domain onto $D$.
The open orbit-map theorem for
$\mathcal K_C\curvearrowright C$ makes every
orbit of this open stabilizer open in $C$.
Since $C$ is connected, the stabilizer acts
transitively on $C$.

There is an accessible point of $\partial D$:
an interior point and a closest boundary point
give a straight end-cut in a coordinate copy of
the plane. Its boundary endpoint belongs to $C$.
Transporting this end-cut by the domain
stabilizer proves that every point of $C$
is accessible from this same $D$.
This holds for every complementary domain.

There are exactly two such domains.
There are at least two because $C$ separates.
For three domains, choose three distinct points
$c_1,c_2,c_3\in C$. Inside each domain, join
end-cuts to these points to an embedded tree
with these three boundary leaves and one
interior branch point. Such a tree is obtained
by joining two end-cuts to an arc and stopping
a third connecting arc at its first intersection.
The three trees have disjoint interiors and
share just their specified boundary leaves.
Their union embeds $K_{3,3}$ in the plane,
contradicting the bipartite planar Euler
inequality $9\leq2\cdot6-4$.

\emph{The proper line conclusion.}
If $C$ is unbounded, put
$C^*=C\cup\{\infty\}$ in the sphere; if $C$
is bounded, put $C^*=C$.
This is a compact connected proper subset
with exactly two complementary domains.
Every point except possibly the added
point $\infty$ is accessible from both.
Choose a stereographic pole in one complementary
domain. Deleting that pole leaves that domain
connected, and accessibility at its boundary
is unaffected, since end-cuts can be chosen
to avoid the pole. In the resulting plane
both inaccessible subsets of $C^*$ have
at most one point.
Polulyakh's converse Jordan theorem
\cite[Theorem 1.1 and Corollary 1.1]{C_Polulyakh2000}
therefore makes $C^*$ a Jordan curve.
The bounded alternative is impossible:
the real-valued open map $F$ is constant
on that Jordan curve, and its restriction
to the bounded disk either has an interior
extremum or is constant on a nonempty
open set. Both contradict openness.
Thus $C$ is a Jordan curve with the
compactification point removed, a properly
embedded line. Transitivity of $\mathcal K$
transfers this conclusion to every component
of every fiber.

Downstairs, a fiber of $\Phi_z$ is homogeneous
under $L_z$, and each of its component
stabilizers is transitive by the same
orbit argument. A nondegenerate compact
component is therefore homogeneous by
ambient neighborhood homeomorphisms.
Brechner--Erkama
\cite[Theorem 2]{C_BrechnerErkama1979}
makes it a Jordan curve.

\emph{An exact disconnected-fiber example.}
On $r>1$ in the plane put $s=\log(r-1)$.
Let $\beta_b$ be the time-$b$ flow of
$\sin\delta\,\partial_\delta$ on $\R$.
It fixes $\pi\mathbb Z$, is $2\pi$-periodic
as a lifted circle map, and
$|\beta_b(\delta)-\delta|\leq\pi$.
Define
\[
\begin{aligned}
 A_a(r,\theta)&=(r,\theta+a) &&(r\leq1),\\
 A_a(r,\theta)&=(1+e^a(r-1),\theta+a)&&(r>1),\\
 B_b|_{\{r\leq1\}}&=\operatorname{id},\\
 B_b(s,\theta)&=(\theta+\beta_b(s-\theta),\theta)&&(r>1).
\end{aligned}
                                                        \tag{C.42.5}
\]
The last first coordinate is $s$.
Its expression is independent of the chosen
representative of $\theta$.
The bound on $\beta_b-\operatorname{id}$
gives continuity at $r=1$ and properness at
infinity for $B_b$ and $B_{-b}$.
The $A$-formulas have these properties directly.
They are continuous orientation-preserving
one-parameter actions and commute, since
$A_a$ preserves $s-\theta$.

The open spiral strip
\[
 (t,\delta)\longmapsto
  (1+e^{t+\delta},\,t\bmod2\pi),
 \qquad t\in\R,\quad0<\delta<\pi,
                                                        \tag{C.42.6}
\]
is an open orbit $O$. This parametrization is
injective because an angular ambiguity is
$2\pi k$, while two $\delta$-values differ
by less than $\pi$.
In these coordinates $A_a$ translates $t$
and $B_b$ sends $\delta$ to $\beta_b(\delta)$.
The resulting group is $\R^2$, simply
transitive on $O$ and closed in the plane
homeomorphism group. Indeed, the open-orbit
preservation argument at the start of the proof
shows that any homeomorphic limit preserves
$O$, using both the forward and inverse limits.
Its evaluation at $(0,\pi/2)$ therefore remains
inside $O$ and recovers finite parameters
$a=t$ and $b=\log\tan(\delta/2)$ continuously.
The parameter limits recover the limiting
plane homeomorphism, proving closedness.
A triangular parameter loop
$0\to(1,0)\to(0,1)\to0$ generates this
group and contracts by scaling both parameters.

For the actual commuting operator subgroup
$H=\langle A_{-2\pi}\rangle$, its powers
$h_n=A_{-2\pi n}$ converge on compact
subsets of $O$ to
\[
 \Phi(t,\delta)=(1,t\bmod2\pi).
                                                        \tag{C.42.7}
\]
This is exactly the quotient associated to
the limit of $h_n(0,\pi/2)$.
Its fiber over $\theta_0$ is the disjoint
union of the lines $t=\theta_0+2\pi k$,
$k\in\mathbb Z$. Its real lift is $t$.
All maps and the ambient contraction have
been specified, so this example concerns
disconnected whole fibers rather than a
nonzero centralizer fundamental-group class.
The executed symbolic certificate
\path{code/part_b/c42_spiral_fibre_certificate.py}
checks the flow law and differential equation
in the exact half-angle coordinate,
the $A$--$B$ commutation identities, the
integer-labelled whole fibers and their real
lift, and the based triangular parameter loop.
It uses no numerical approximations; the
topological conclusions are proved above.
\end{proof}

\paragraph{Scope of the reduction.}
C.42 proves the component statement for the
specified real-valued lift. It does not
identify a lifted component with a whole
downstairs connected component without
checking the actual covering restriction.
It does not yet give local product charts
for the component decomposition, simultaneous
charts for several boundary quotients, or
the frontier-compatible contraction needed
in C.12. Those global extension questions
remain open.
